\documentclass[UTF8,aspectratio=169]{ctexbeamer}
\usetheme{Madrid}
\usecolortheme{default}
\usepackage{amsmath}
\usepackage{booktabs}
\usepackage{siunitx}
\usepackage{graphicx}
\usepackage{array}
\usepackage{xcolor}
\usepackage{xurl}
\usepackage{colortbl}
\usepackage{bm}

\newcommand{\good}[1]{\textcolor{green!45!black}{#1}}
\newcommand{\warn}[1]{\textcolor{red!65!black}{#1}}
\newcommand{\bluenote}[1]{\textcolor{blue!65!black}{#1}}
\newcommand{\src}[1]{\vfill{\tiny\textcolor{gray!75!black}{Source: \url{#1}}}}
\newcolumntype{L}[1]{>{\raggedright\arraybackslash}p{#1}}
\newcolumntype{R}[1]{>{\raggedleft\arraybackslash}p{#1}}
\newcolumntype{C}[1]{>{\centering\arraybackslash}p{#1}}

\definecolor{winrow}{RGB}{220,240,220}
\definecolor{warnrow}{RGB}{255,235,230}

\title[Stage-2 Anchor Ablation]{Stage-2 Anchor Ablation\\LONO Validation of the Early-Time Physics Prior}
\author{Mie Spray Postprocessing Project}
\date{\today}

\begin{document}

\begin{frame}
  \titlepage
\end{frame}

% ──────────────────────────────────────────────────────────────
\begin{frame}{Message}
\small
\textbf{Claim:} adding a soft $\bm{\mu(t{\to}0){\approx}0}$ anchor to the Stage-2 NLL loss
improves out-of-nozzle generalisation. Extending the anchor to also constrain $\sigma$
brings no further gain (and slightly over-tightens uncertainty).

\vspace{0.5em}
\textbf{This deck:}
\begin{itemize}
  \item recaps the physical motivation: penetration must start at $0$ at injection onset,
        so $\mu$ and $\sigma$ near $t{=}0$ should be small
  \item defines the three anchor ablations: \texttt{no\_anchor}, \texttt{mu\_anchor},
        \texttt{mu\_sigma\_anchor}
  \item shows per-fold and mean post-Stage-3 RMSE/MAE/coverage for all three under 5-fold LONO
  \item confirms (i) \texttt{mu\_anchor} wins on RMSE, MAE, and 2$\sigma$ coverage,
        (ii) \texttt{no\_anchor} catastrophically fails on the hardest non-Nozzle0 fold,
        (iii) adding $\sigma$-anchor narrows predicted $\sigma$ but does not improve RMSE
\end{itemize}

\vspace{0.5em}
\textbf{Validated pipeline:} \texttt{run\_stage2\_lono\_ablation.py}, seed 42, 5 folds $\times$ 3 ablations
(uncensored eval points per fold $\in\{311149,\,60369,\,33017,\,28139,\,25705\}$);
Stage-1 backbone fixed to \texttt{a\_dp050\_plus\_pressures} (Stage-1 winner).
\end{frame}

% ──────────────────────────────────────────────────────────────
\begin{frame}{Background: The Early-Time Physics Anchor}
\small
\textbf{Physical constraint at injection onset:}
\[
  S(t{=}0) = 0,\qquad \sigma\bigl(S(t{=}0)\bigr) \;\text{should be small.}
\]
Without an explicit penalty the NLL trains the network to optimise the \emph{likelihood of the
observed points}, which contain almost no samples at $t < \SI{0.1}{ms}$ (laser onset truncation,
shadow latency, etc.). The network therefore drifts to large $|\mu|$ and inflated $\sigma$
near $t{=}0$ just because that region is data-poor.

\vspace{0.5em}
\textbf{Anchor penalty} (added to the NLL):
\[
  \mathcal{R}_{\mu} = \frac{\sum_t w(t)\,\mu(t)^2}{\sum_t w(t)},\qquad
  \mathcal{R}_{\sigma} = \frac{\sum_t w(t)\,\bigl[\max(0,\,\sigma(t)-\sigma_0)\bigr]^2}{\sum_t w(t)},
\]
\[
  w(t) = \max\!\Bigl(0,\;1 - \tfrac{t}{\tau}\Bigr),\qquad
  \tau = \text{anchor\_window\_ms (early-time ramp)}.
\]
$w(t)$ is a triangular kernel: weight $1$ at $t{=}0$, decaying linearly to $0$ at $t{=}\tau$.

\vspace{0.5em}
\textbf{Final Stage-2 loss:}
\[
  \mathcal{L}_2 \;=\; \mathrm{NLL}\bigl(y\,;\,\mu(x),\sigma(x)\bigr)
              \;+\; \lambda_\mu\,\mathcal{R}_\mu
              \;+\; \lambda_\sigma\,\mathcal{R}_\sigma.
\]
\end{frame}

% ──────────────────────────────────────────────────────────────
\begin{frame}{Variant Design: What Each Configuration Tests}
\small
\begin{center}
\renewcommand{\arraystretch}{1.25}
\begin{tabular}{@{}L{3.4cm} C{1.2cm} C{1.2cm} L{6.5cm}@{}}
\toprule
\textbf{Variant} & $\bm{\lambda_\mu}$ & $\bm{\lambda_\sigma}$ & \textbf{What it tests} \\
\midrule
\texttt{no\_anchor}        & 0       & 0       & pure NLL baseline; no physics prior at $t{=}0$ \\
\texttt{mu\_anchor}        & \good{$10^{-2}$} & 0       & \good{soft anchor only on $\mu$ — penalises non-zero penetration at injection onset} \\
\texttt{mu\_sigma\_anchor} & $10^{-2}$ & $10^{-3}$ & adds soft anchor on excess-$\sigma$ above the floor $\sigma_0$ \\
\bottomrule
\end{tabular}
\end{center}

\vspace{0.6em}
\textbf{Identical across ablations:}
\begin{itemize}\setlength\itemsep{2pt}
  \item Stage-1 backbone: \texttt{a\_dp050\_plus\_pressures} (locked from Stage-1 LONO winner)
  \item Stage-3 chain: \texttt{anchor\_off} distillation (identical recipe)
  \item Same train/val split, same seed (42), same nozzle held out per fold
\end{itemize}

\vspace{0.4em}
\textbf{Per-fold pipeline:} Stage-1 (1 run) $\to$ Stage-2 (3 ablations) $\to$ Stage-3 anchor\_off (3 chains) $\to$
post-Stage-3 RMSE eval on the held-out nozzle's uncensored points.

\vspace{0.3em}
\footnotesize
Total trainings: $5 \times (1{+}3{+}3) = 35$ runs; reported metrics are the post-Stage-3 eval (end-task quality).
\end{frame}

% ──────────────────────────────────────────────────────────────
\begin{frame}{LONO Results: Summary (Mean $\pm$ Std Across 5 Folds)}
\small
\begin{center}
\renewcommand{\arraystretch}{1.22}
\setlength{\tabcolsep}{4pt}
\begin{tabular}{@{}lrrrrr@{}}
\toprule
\textbf{Ablation} & \textbf{RMSE (mm)} & \textbf{MAE (mm)} & \textbf{Bias (mm)} & \textbf{cov$_{2\sigma}$} & \textbf{$\bar\sigma$ (mm)} \\
\midrule
\rowcolor{winrow}
\texttt{mu\_anchor}        & \textbf{12.73 $\pm$ 12.44} & \textbf{9.86 $\pm$ 10.22} & $-7.11 \pm 11.51$ & \textbf{0.775 $\pm$ 0.34} & 5.68 $\pm$ 1.44 \\
\texttt{mu\_sigma\_anchor} & 13.05 $\pm$ 12.46 & 10.06 $\pm$ 10.33 & $-7.59 \pm 11.45$ & 0.764 $\pm$ 0.35 & \textbf{5.49 $\pm$ 1.55} \\
\rowcolor{warnrow}
\texttt{no\_anchor}        & 14.54 $\pm$ 12.36 & 10.80 $\pm$ 9.96 & $-8.50 \pm 11.03$ & 0.763 $\pm$ 0.34 & 6.11 $\pm$ 1.58 \\
\bottomrule
\end{tabular}
\end{center}

\vspace{0.5em}
\textbf{Headline numbers, excluding the outlier Nozzle0 fold (4-fold mean RMSE):}
\begin{center}
\renewcommand{\arraystretch}{1.2}
\begin{tabular}{@{}lc@{}}
\toprule
\textbf{Ablation} & \textbf{RMSE w/o Nozzle0 (mm)} \\
\midrule
\rowcolor{winrow}
\texttt{mu\_anchor}        & \textbf{7.18} \\
\texttt{mu\_sigma\_anchor} & 7.53 \\
\rowcolor{warnrow}
\texttt{no\_anchor}        & 9.61 \\
\bottomrule
\end{tabular}
\end{center}

\vspace{0.3em}
\footnotesize
\good{Green}: anchored variants. \warn{Red}: no-anchor baseline.
\texttt{mu\_anchor} wins on RMSE, MAE, and 2$\sigma$-coverage; \texttt{mu\_sigma\_anchor}
has the tightest predicted $\sigma$ (5.49 mm) but loses 0.32 mm RMSE — over-constraining the
uncertainty channel hurts the mean.
\end{frame}

% ──────────────────────────────────────────────────────────────
\begin{frame}{LONO Results: Per-Fold Breakdown}
\footnotesize
\begin{center}
\renewcommand{\arraystretch}{1.08}
\setlength{\tabcolsep}{4pt}
\begin{tabular}{@{}lrrrrr|r@{}}
\toprule
\textbf{Ablation} &
\textbf{Nozzle0} & \textbf{Nozzle2} & \textbf{Nozzle5} & \textbf{Nozzle4} & \textbf{Nozzle1} &
\textbf{Mean} \\
& \scriptsize(311149) & \scriptsize(60369) & \scriptsize(33017) & \scriptsize(28139) & \scriptsize(25705) & \\
\midrule
\multicolumn{7}{l}{\textit{RMSE (mm)}} \\
\rowcolor{winrow}
\texttt{mu\_anchor}        & 34.92 & \good{7.76} & 5.60 & 8.20 & 7.16 & \textbf{12.73} \\
\texttt{mu\_sigma\_anchor} & 35.12 & 10.04       & 5.27 & 7.73 & 7.10 & 13.05 \\
\rowcolor{warnrow}
\texttt{no\_anchor}        & 34.25 & \warn{19.18} & 5.11 & 7.57 & 6.56 & 14.54 \\
\midrule
\multicolumn{7}{l}{\textit{MAE (mm)}} \\
\texttt{mu\_anchor}        & 28.11 & 5.72  & 4.26 & 5.84 & 5.37 & 9.86 \\
\texttt{mu\_sigma\_anchor} & 28.44 & 6.91  & 3.89 & 5.79 & 5.29 & 10.06 \\
\texttt{no\_anchor}        & 27.53 & 12.46 & 3.79 & 5.41 & 4.82 & 10.80 \\
\midrule
\multicolumn{7}{l}{\textit{Coverage at $2\sigma$ (target 0.95)}} \\
\texttt{mu\_anchor}        & .171 & .856 & .975 & .912 & .963 & .775 \\
\texttt{mu\_sigma\_anchor} & .149 & .801 & .981 & .929 & .959 & .764 \\
\texttt{no\_anchor}        & .191 & .690 & .985 & .965 & .982 & .763 \\
\bottomrule
\end{tabular}
\end{center}

\vspace{0.1em}
\scriptsize
Post-Stage-3 metrics on held-out nozzle's uncensored CDF P50 points. Number under each
nozzle = eval-point count. Columns sorted by descending point count; Nozzle0 dominates the
mean by sample weight.
\end{frame}

% ──────────────────────────────────────────────────────────────
\begin{frame}{Discussion: Three Consistent Patterns}
\small
\begin{block}{1. \good{$\mu$-anchor stabilises the hardest non-outlier fold}}
On Nozzle2 (the second-largest fold by sample count and the only fold where the model is asked
to extrapolate to a new injector geometry within the modern family), \texttt{no\_anchor}
collapses to \warn{19.18\,mm RMSE} while \texttt{mu\_anchor} holds at 7.76\,mm. The anchor
prevents the NLL from fitting an arbitrary near-$t{=}0$ baseline using the data-poor early
window as free parameters.
\end{block}

\begin{block}{2. \warn{Adding $\sigma$-anchor doesn't recover the lost RMSE}}
\texttt{mu\_sigma\_anchor} achieves the tightest predicted $\sigma$ (5.49\,mm mean) but its
RMSE is 0.32\,mm worse than \texttt{mu\_anchor}. Pushing $\sigma$ down at $t{=}0$ propagates
through Stage-3 distillation as a more confident-but-biased $\mu$ prediction, hurting accuracy
without improving calibration (cov$_{2\sigma}$ 0.764 vs 0.775).
\end{block}

\begin{block}{3. \bluenote{Easy folds (Nozzle1/4/5) are insensitive to the anchor choice}}
On the three modern-design nozzles with $\sim$25--33k eval points, all three ablations sit in
a narrow 5.1--8.2\,mm RMSE band. The anchor only matters when the LONO fold creates a
genuine OOD challenge (Nozzle0 for geometric reasons, Nozzle2 for sample-size reasons).
\end{block}
\end{frame}

% ──────────────────────────────────────────────────────────────
\begin{frame}{Discussion: Nozzle0 is Still the Hardest Fold}
\small
\begin{columns}[T]
\column{0.55\textwidth}
Nozzle0 (\texttt{BC20220627\_HZ\_Nozzle0}) carries \textbf{311{,}149 uncensored eval points}
— about 5$\times$ more than the next-largest fold (Nozzle2: 60{,}369). It is also the
\emph{only} first-generation nozzle in the dataset.

\vspace{0.5em}
All three ablations land in the same catastrophic band on this fold:
\begin{itemize}\setlength\itemsep{2pt}
  \item RMSE 34--35\,mm
  \item bias $\approx -27$\,mm (systematic under-prediction)
  \item 2$\sigma$ coverage 0.15--0.19
\end{itemize}

\vspace{0.5em}
\textbf{Interpretation:} the Nozzle0 failure is not an anchor problem; it is a fundamental
\emph{out-of-design-family} extrapolation. None of the in-distribution Stage-2 regularisers
can compensate for a missing geometric class. This was already known from the Stage-1
ablation and is reproduced here at the Stage-3 output level.

\column{0.42\textwidth}
\begin{center}
\renewcommand{\arraystretch}{1.2}
\small
\begin{tabular}{@{}lrr@{}}
\toprule
\textbf{Fold} & \textbf{Points} & \textbf{Best RMSE} \\
\midrule
Nozzle0 & 311{,}149 & 34.25\,mm \\
Nozzle2 & 60{,}369  & 7.76\,mm \\
Nozzle5 & 33{,}017  & 5.11\,mm \\
Nozzle4 & 28{,}139  & 7.57\,mm \\
Nozzle1 & 25{,}705  & 6.56\,mm \\
\bottomrule
\end{tabular}
\end{center}
\vspace{0.5em}
\footnotesize
Best RMSE = minimum across the three ablations for that fold.
Excluding Nozzle0 collapses the cross-ablation gap to 7.18\,mm
(\texttt{mu\_anchor}) vs 9.61\,mm (\texttt{no\_anchor}), a 2.4\,mm OOD margin from the
anchor alone.
\end{columns}
\end{frame}

% ──────────────────────────────────────────────────────────────
\begin{frame}{Conclusion: $\mu$-Anchor is the Stage-2 LONO Winner}
\small
\begin{block}{What the ablation proves}
\begin{enumerate}
  \item \textbf{The early-time $\mu$-anchor is load-bearing.} It buys 1.8\,mm mean RMSE
        (14.54 $\to$ 12.73) over no-anchor with Nozzle0 included, and 2.4\,mm
        (9.61 $\to$ 7.18) excluding Nozzle0.
  \item \textbf{The anchor's value shows up on hard OOD folds, not easy ones.} On Nozzle1/4/5
        all variants are within 0.5\,mm; on Nozzle2 the anchor saves 11+\,mm. The
        regularisation only matters when the model is forced to extrapolate.
  \item \textbf{$\sigma$-anchor is not additive.} Tightening $\sigma$ at $t{=}0$ trades
        $\bar\sigma$ (5.68 $\to$ 5.49\,mm) for a worse RMSE (12.73 $\to$ 13.05\,mm) and
        slightly worse calibration. \texttt{mu\_anchor} is the Pareto winner.
  \item \textbf{Production variant for downstream Stage-3 LONO and multi-seed retrains:}
        \texttt{mu\_anchor} ($\lambda_\mu{=}10^{-2}$, $\lambda_\sigma{=}0$).
\end{enumerate}
\end{block}

\vspace{0.3em}
\textbf{Experiment:} \texttt{run\_stage2\_lono\_ablation.py} --- 3 ablations $\times$ 5 LONO
folds, seed 42, Stage-1 = \texttt{a\_dp050\_plus\_pressures}, Stage-3 = \texttt{anchor\_off}.\\
\textbf{Output:} \texttt{MLP/runs\_mlp/stage2\_lono\_ablation\_20260520\_235200/}
\end{frame}

\end{document}
